\section*{Розділ VI. Студентські ради гуртожитків}
\addcontentsline{toc}{section}{Розділ VI. Студентські ради гуртожитків}

\subsection*{6.1. Статус та завдання СРГ}
\addcontentsline{toc}{subsection}{6.1. Статус та завдання СРГ}
    6.1.1. Студентська рада гуртожитку (далі – СРГ) є основним органом студентського самоврядування на рівні гуртожитку, який діє відповідно до Закону України "Про вищу освіту", Статуту Університету, Положення про ОСС та цього Положення.

    6.1.2. СРГ є підзвітною та підконтрольною мешканцям гуртожитку (через процедуру виборів та звітування Голови), Конференції студентів Університету (КСУ) та Студентській раді студмістечка в межах їхніх повноважень.

    6.1.3. Основним завданням СРГ є представництво та захист прав та інтересів мешканців свого гуртожитку, а також реалізація завдань студентського самоврядування у побутовій сфері, визначених Положенням про ОСС та цим Положенням, на рівні гуртожитку.

\subsection*{6.2. Формування та склад СРГ}
\addcontentsline{toc}{subsection}{6.2. Формування та склад СРГ}
    6.2.1. СРГ очолює Голова, який обирається мешканцями гуртожитку шляхом прямого таємного голосування строком на 1 (один) рік відповідно до процедури, визначеної Положенням про ЦВКс. Одна і та сама особа не може обіймати посаду Голови СРГ більше двох строків поспіль.

    6.2.2. Інші члени СРГ залучаються до складу Головою СРГ на підставі особистої заяви студента, який є мешканцем відповідного гуртожитку.

    6.2.3. Голова СРГ видає розпорядження про включення студента до складу СРГ, якщо відсутні підстави для відмови, визначені у п. 6.2.4.

    6.2.4. Рішення про відмову студенту у включенні до складу СРГ приймається \textbf{колегіально СРГ} більшістю голосів від загального складу виключно з таких підстав:

        \begin{enumerate}[label=\alph*)]
            \item Невідповідність кандидата вимогам (не є мешканцем відповідного гуртожитку);
            \item Попереднє виключення кандидата зі складу ОСС за грубе порушення нормативних документів ОСС або вчинення дій, що завдали значної шкоди репутації ОСС.
        \end{enumerate}

    6.2.5. Рішення СРГ про відмову у включенні до складу може бути оскаржене студентом до Студентської уповноваженої делегації (СУД) протягом 7 (семи) календарних днів з моменту його прийняття.

    6.2.6. Загальна кількість членів СРГ (включаючи Голову) з правом голосу обмежується \textbf{8 (вісьмома) особами}.

    6.2.7. Студенти, залучені до складу СРГ понад встановлений ліміт (8 осіб), беруть участь у роботі з правом \textbf{дорадчого голосу} та не можуть претендувати на додаткові бали чи інші заохочення, передбачені для членів ОСС з правом голосу.

    6.2.8. СРГ має право залучати до своєї роботи інших студентів Університету денної форми навчання, які не є мешканцями відповідного гуртожитку. Такі залучені студенти беруть участь у роботі СРГ з правом дорадчого голосу.

    6.2.9. Секретар СРГ веде актуальний реєстр членів СРГ із зазначенням дати включення до складу та чітким розмежуванням членів з правом голосу та членів з правом дорадчого голосу.

    6.2.10. У виняткових випадках, за наявності обґрунтованої потреби, Конференція студентів Університету (КСУ) за поданням Студентської ради студмістечка (СР СМ) або Студентської уповноваженої делегації (СУД) може прийняти рішення про збільшення максимальної кількості членів відповідної СРГ з правом голосу понад ліміт, встановлений п. 6.2.6 цього Положення. Таке рішення має визначати новий ліміт та, за потреби, термін його дії.

\subsection*{6.3. Структура СРГ}
\addcontentsline{toc}{subsection}{6.3. Структура СРГ}
    6.3.1. Обов'язковими посадами в СРГ є:

        \begin{enumerate}[label=\alph*)]
            \item Голова СРГ;
            \item Заступник (або заступники) Голови СРГ — виконує обов'язки Голови за його відсутності або дорученням, координує роботу визначених напрямів;
            \item Секретар СРГ — відповідає за ведення протоколів засідань, діловодство, ведення реєстру членів та архіву СРГ.
        \end{enumerate}

    6.3.2. Заступник(-и) Голови та Секретар СРГ призначаються Головою СРГ з числа членів СРГ з правом голосу.

    6.3.3. СРГ за поданням Голови може створювати постійні або тимчасові структурні підрозділи (відділи, комітети, проєктні групи) для виконання специфічних завдань. Рішення про створення або ліквідацію таких підрозділів приймається колегіально СРГ.

    6.3.4. Керівники структурних підрозділів призначаються Головою СРГ з числа членів СРГ (переважно з правом голосу).

\subsection*{6.4. Повноваження Голови СРГ}
\addcontentsline{toc}{subsection}{6.4. Повноваження Голови СРГ}
    6.4.1. Голова СРГ, окрім повноважень, визначених у п. 6.1.3 та інших розділах цього Положення:

        \begin{enumerate}[label=\alph*)]
            \item Здійснює загальне керівництво діяльністю СРГ.
            \item Забезпечує формування складу СРГ шляхом розгляду заяв студентів та видання розпоряджень про включення до складу відповідно до п. 6.2.2-6.2.3.
            \item Ініціює перед СРГ розгляд питання про відмову у включенні до складу відповідно до п. 6.2.4.
            \item Призначає Заступника(-ів) Голови, Секретаря та керівників структурних підрозділів СРГ.
            \item Організовує та контролює роботу членів та структурних підрозділів СРГ.
            \item Скликає та веде засідання СРГ.
            \item Видає розпорядження організаційно-виконавчого характеру в межах своєї компетенції, що не належать до виключної компетенції колегіального органу СРГ.
            \item Вносить на розгляд СРГ подання про дострокове припинення повноважень (виключення) члена СРГ з підстав, визначених цим Положенням.
            \item Представляє СРГ у відносинах з адміністрацією гуртожитку та студмістечка, СР СМ, іншими органами та організаціями.
            \item Підписує протоколи засідань, рішення та інші офіційні документи СРГ.
            \item Звітує про свою діяльність та діяльність СРГ перед мешканцями гуртожитку та СР СМ.
        \end{enumerate}

\subsection*{6.5. Припинення повноважень членів СРГ}
\addcontentsline{toc}{subsection}{6.5. Припинення повноважень членів СРГ}
    6.5.1. Повноваження члена СРГ (включаючи Голову) припиняються у зв'язку із закінченням строку, на який його було обрано/призначено (1 рік).

    6.5.2. Повноваження члена СРГ (включаючи Голову) припиняються достроково та автоматично у разі:

        \begin{enumerate}[label=\alph*)]
            \item Подання особистої заяви про складання повноважень;
            \item Втрати статусу студента Університету;
            \item Втрати статусу мешканця відповідного гуртожитку (виселення, переселення до іншого гуртожитку);
            \item Набрання законної сили обвинувальним вироком суду щодо нього;
            \item Смерті або визнання його безвісно відсутнім чи померлим.
        \end{enumerate}

    6.5.3. Повноваження \textbf{Голови} СРГ можуть бути достроково припинені (він може бути відкликаний/усунутий з посади) також за рішенням:

        \begin{enumerate}[label=\alph*)]
            \item Мешканців гуртожитку (відкликання виборцями) у порядку, визначеному Положенням про ЦВКс;
            \item Конференції студентів Університету (КСУ) за висловленням недовіри або з інших підстав, передбачених Положенням про ОСС, у порядку, визначеному Регламентом КСУ та/або Положенням про ОСС;
            \item Студентської уповноваженої делегації (СУД) у випадках та порядку, передбачених Положенням про СУД.
        \end{enumerate}

    6.5.4. Повноваження члена СРГ (\textbf{окрім Голови}) можуть бути достроково припинені за рішенням СРГ (прийнятим колегіально більшістю голосів від загального складу) за поданням Голови СРГ у разі:

        \begin{enumerate}[label=\alph*)]
            \item Систематичного невиконання обов'язків члена СРГ, зокрема неучасті у засіданнях чи роботі структурного підрозділу;
            \item Грубого порушення нормативних документів ОСС або Статуту Університету;
            \item Вчинення дій, що завдали значної шкоди репутації ОСС Університету.
        \end{enumerate}

    6.5.5. Перед прийняттям рішення відповідно до п. 6.5.4 СРГ зобов'язана надати можливість члену, стосовно якого розглядається питання, надати пояснення.

    6.5.6. Рішення СРГ про дострокове припинення повноважень члена (окрім Голови) може бути оскаржене до Студентської уповноваженої делегації (СУД) протягом 7 (семи) календарних днів з моменту його прийняття.

    \subsubsection*{6.5.7. Призначення виконувача обов'язків Голови СРГ}
        6.5.7.1. У разі дострокового припинення повноважень Голови СРГ, його тимчасової неможливості виконувати свої обов'язки (через хворобу, відрядження тощо) або інших обставин, що унеможливлюють виконання ним своїх повноважень, призначається виконувач обов'язків (в.о.) Голови СРГ.

        6.5.7.2. В.о. Голови СРГ може бути призначений колегіальним рішенням СРГ на термін не більше 1 (одного) місяця. Для призначення в.о. на більший термін необхідне рішення КСУ або СУД за поданням СР СМ.

        6.5.7.3. В.о. Голови СРГ має повний обсяг повноважень Голови СРГ. Призначення на посаду в.о. Голови СРГ не надає права голосу в СР СМ та не змінює наявного статусу щодо права голосу в СРГ. Якщо в.о. не є членом СРГ з правом голосу, він не набуває такого права у зв'язку з призначенням. Якщо рішення про призначення в.о. прийнято КСУ або СУД, у ньому можуть зазначатися додаткові обмеження повноважень.

        6.5.7.4. Передача справ в.о. Голови СРГ здійснюється відповідно до процедури, визначеної розділом 4.6 Положення про ОСС.

    \subsubsection*{6.5.8. Призначення виконувачів обов'язків інших посадових осіб СРГ}
        6.5.8.1. У разі дострокового припинення повноважень, тимчасової неможливості виконувати обов'язки або інших обставин, що унеможливлюють виконання повноважень Заступника(-ів) Голови, Секретаря СРГ або керівників структурних підрозділів, призначення виконувача обов'язків здійснюється в тому ж порядку, що й призначення відповідних посадових осіб.

        6.5.8.2. Термін повноважень, обсяг повноважень та порядок передачі справ для таких в.о. визначаються рішенням про їх призначення.

\subsection*{6.6. Основні повноваження та порядок роботи СРГ}
\addcontentsline{toc}{subsection}{6.6. Основні повноваження та порядок роботи СРГ}
    6.6.1. СРГ виконує завдання та повноваження, визначені у п. 6.1.3. Основні напрями діяльності включають, але не обмежуються:

        \begin{enumerate}[label=\alph*)]
            \item Представництво та захист прав мешканців гуртожитку перед адміністрацією гуртожитку, студмістечка та Університету.
            \item Сприяння поліпшенню умов проживання, побуту та відпочинку мешканців.
            \item Участь у процесах поселення, переселення та виселення в межах гуртожитку.
            \item Сприяння дотриманню мешканцями правил внутрішнього розпорядку гуртожитку.
            \item Організація та проведення заходів для мешканців гуртожитку.
            \item Інформаційне забезпечення мешканців гуртожитку.
            \item Взаємодія з СР СМ та іншими ОСС.
        \end{enumerate}

    6.6.2. Основною формою роботи СРГ є засідання, які проводяться за потребою, але не рідше одного разу на місяць протягом навчального семестру.

    6.6.3. У засіданні СРГ беруть участь з правом голосу члени СРГ в межах встановленої квоти (до 8 осіб) та з правом дорадчого голосу — інші залучені члени. Засідання є правомочним, якщо на ньому присутні більше половини від складу членів з правом голосу.

    6.6.4. Рішення СРГ приймаються більшістю голосів від загального складу членів з правом голосу, якщо інше не передбачено цим Положенням. Процедурні питання можуть вирішуватися більшістю від присутніх членів з правом голосу.

    6.6.5. \textbf{Виключно колегіально} СРГ приймає рішення з таких питань:

        \begin{enumerate}[label=\alph*)]
            \item Відмова у включенні студента до складу СРГ (п. 6.2.4);
            \item Дострокове припинення повноважень (виключення) члена СРГ (п. 6.5.4);
            \item Створення або ліквідація структурних підрозділів СРГ;
            \item Затвердження плану роботи та звіту про діяльність СРГ;
            \item Затвердження позиції СРГ щодо погодження рішень адміністрації гуртожитку та студмістечка з питань, що належать до компетенції ОСС (напр., щодо поселення, виселення, правил внутрішнього розпорядку);
            \item Інші питання, віднесені до виключної компетенції СРГ цим Положенням або за рішенням самої СРГ.
        \end{enumerate}

    6.6.6. Хід засідання фіксується у протоколі, який веде Секретар СРГ.

    6.6.7. СРГ здійснює свою діяльність відповідно до плану роботи, затвердженого колегіально на початку навчального року.

    6.6.8. Внутрішні спори та суперечності між членами СРГ або між Головою та членами СРГ вирішуються шляхом обговорення та прийняття колегіального рішення на засіданні СРГ. У разі неможливості дійти згоди, будь-яка зі сторін може звернутися до СР СМ для надання допомоги у вирішенні спору (медіації).

    6.6.9. Рішення про зміну статусу члена СРГ щодо права голосу (надання або позбавлення права голосу) в межах встановленого п. 6.2.6 ліміту приймається колегіально СРГ кваліфікованою більшістю голосів — не менше 2/3 (двох третин) від загального складу членів з правом голосу. Таке рішення не може призводити до перевищення ліміту у 8 членів з правом голосу. Секретар вносить відповідні зміни до реєстру членів СРГ.

\subsection*{6.7. Взаємодія СРГ з СР СМ}
\addcontentsline{toc}{subsection}{6.7. Взаємодія СРГ з СР СМ}
    6.7.1. СРГ зобов'язані брати участь у координаційних заходах, що проводяться СР СМ.

    6.7.2. СРГ надають СР СМ інформацію та звіти про свою діяльність у порядку та строки, встановлені СР СМ.

    6.7.3. СРГ розглядають рекомендації СР СМ щодо організації своєї роботи та проведення заходів та інформують СР СМ про результати розгляду.

    6.7.4. СРГ сприяють виконанню рішень СР СМ, що стосуються загальних питань студмістечка, на рівні гуртожитку.

\subsection*{6.8. Реорганізація гуртожитку}
\addcontentsline{toc}{subsection}{6.8. Реорганізація гуртожитку}
    6.8.1. У разі закриття, реорганізації або тимчасового виведення гуртожитку з експлуатації, порядок припинення або тимчасового зупинення діяльності відповідної СРГ, передачі документації та активів, а також захисту інтересів мешканців, переселених до інших гуртожитків, визначається рішенням СР СМ відповідно до загальних засад, визначених Положенням про ОСС.

    \subsubsection*{6.8.2. Тимчасове управління від СР СМ}

        6.8.2.1. У випадку, якщо проведення виборів Голови СРГ є неможливим у встановлені терміни, або якщо діяльність СРГ фактично припинена з інших причин (відсутність кворуму, масове складання повноважень членами тощо), КСУ за поданням СР СМ або СУД може прийняти рішення про запровадження тимчасового управління справами студентського самоврядування відповідного гуртожитку.

        6.8.2.2. Тимчасове управління здійснюється Тимчасовим комітетом, що формується та призначається СР СМ.

        6.8.2.3. Повноваження Тимчасового комітету за замовчуванням обмежуються підтриманням базової операційної діяльності, забезпеченням мінімально необхідного представництва інтересів мешканців гуртожитку та першочерговою організацією та сприянням проведенню виборів Голови СРГ. За окремим рішенням КСУ Тимчасовому комітету можуть бути надані розширені повноваження.

        6.8.2.4. Тимчасове управління триває до моменту обрання та початку роботи нового Голови СРГ та формування ним основного складу СРГ, але не може перевищувати 3 (трьох) місяців. Якщо протягом цього терміну вибори не проведено, питання про подальші дії виноситься на розгляд КСУ.

        6.8.2.5. Тимчасовий комітет є підзвітним СР СМ та КСУ.
